% ===================================================================
\section{Apertura y problema}

\diapo{Portada}{PULSO}{0:30}{0:30}{Quién presenta y de qué trata.}
\begin{decir}
Buenos días. Mi nombre es Juan Diego Martínez y presento PULSO: Pagos, Umbrales,
Liquidez y Sincronización Operativa. Es el proyecto del Semillero de Modelamiento de
Sistemas Complejos y Física Computacional, dirigido por el profesor Jhon Hernán Díaz.

En una frase: estudiamos si el momento en que llegan los pagos, y no solo su monto,
ayuda a anticipar quién va a dejar de pagar un crédito.
\end{decir}

\diapo{1}{Los pagos forman una red y los problemas viajan por ella}{1:00}{1:30}{Un
problema de liquidez no se queda donde empieza.}
\begin{decir}
Empiezo con una idea simple. Cada persona y cada empresa paga con el dinero que recibe
de otras. Por eso los pagos forman una red: cada punto es un agente y cada flecha es un
pago.

En el dibujo, A deja de pagar. B y C reciben menos dinero. Si no tienen reservas,
tampoco pueden pagarle a D. Así el problema de A llega a D, aunque D nunca tuvo relación
con A.

En Colombia, Bre-B permite pagos inmediatos entre todos los bancos. Eso es bueno para los
usuarios, pero también hace que un problema de liquidez viaje más rápido.
\transicion{Pero en una red de pagos no solo importa cuánto dinero llega.}
\end{decir}
\begin{terminos}
\begin{description}
\item[Liquidez.] Tener dinero disponible en el momento de pagar. Es distinto de la
solvencia, que compara todo lo que se tiene con todo lo que se debe.
\item[Red de pagos.] Conjunto de agentes (nodos) unidos por flechas de pago con
dirección y monto.
\item[Bre-B.] Sistema de pagos inmediatos del Banco de la República de Colombia. Aquí es
solo motivación: no usamos sus datos.
\end{description}
\end{terminos}

\diapo{2}{No solo importa cuánto se recibe, también cuándo}{1:30}{3:00}{El ingreso
desfasado genera estrés aunque el monto no cambie.}
\begin{decir}
Esta es la idea central del proyecto.

Las dos líneas son el mismo mes de la misma persona. Gana el mismo salario y tiene la
misma cuota, que vence el día 35. En la línea de arriba, el salario llega el
día 28: cuando vence la cuota, ya tiene el dinero. En la línea de abajo, el salario llega
el día 38: durante tres días debe la cuota y no tiene con qué pagarla.

El monto es idéntico. Lo único que cambió es la fecha. A esa diferencia entre cuándo
llega el ingreso y cuándo vencen las obligaciones la llamamos desfase o
desincronización.

Si la persona no tiene un ahorro que cubra esos días, entra en mora. Y si el desfase se
repite cada mes, el problema se acumula. Lo que planteamos es que este efecto, que no
depende del monto, se puede medir y sirve para predecir.
\transicion{Esto lleva directamente a la pregunta del proyecto.}
\end{decir}
\begin{terminos}
\begin{description}
\item[Ciclo de pagos.] Periodo que se repite entre dos pagos del mismo tipo, por
ejemplo un mes de nómina.
\item[Desfase o desincronización.] Diferencia de tiempo entre la llegada del ingreso y
el vencimiento de las obligaciones.
\item[Mora.] Retraso en el pago de una obligación.
\end{description}
\end{terminos}

\diapo{3}{Pregunta de investigación}{1:15}{4:15}{Comparamos contra la red estática
porque eso ya existe en la literatura.}
\begin{decir}
Los bancos deciden a quién prestar con un score: un modelo que estima la probabilidad
de impago con datos del propio cliente, como su ingreso, su deuda y su historial de
mora.

Nuestra pregunta es si la sincronización de los ciclos de pago, medida solo con las
transacciones que el banco ya tiene, mejora esa predicción. Y la comparamos contra dos
referencias: el score clásico, y lo que llamamos red estática, que es saber quién le paga
a quién y cuánto. Ya hay trabajos que usan la red estática; lo nuevo es medir el ritmo.

La hipótesis se puede refutar. Si el desfase genera déficit, la sincronización debe
mejorar la predicción. Si no lo genera, no debe mejorarla. Por eso incluimos un control:
un mundo idéntico donde el desfase no hace daño.
\transicion{Para responder, construimos el modelo en seis pasos.}
\end{decir}
\begin{terminos}
\begin{description}
\item[Score crediticio.] Modelo estadístico, normalmente una regresión logística, que
asigna a cada cliente una probabilidad de impago.
\item[Impago.] Dejar de pagar una obligación. La probabilidad de impago a un año se
llama PD.
\item[Red estática.] Variables sobre quién le paga al cliente y cuánto, sin mirar el
momento de los pagos.
\item[Hipótesis falsable.] Afirmación que predice un resultado que podría no ocurrir.
\item[Control nulo.] Experimento idéntico con el mecanismo apagado. Mide cuánto
``efecto'' produce el método por sí solo.
\end{description}
\end{terminos}
\begin{preguntas}
\pr{¿No es lo mismo que el historial crediticio?}{No. El historial describe al cliente
solo. La red estática describe sus relaciones, y ya aparece en la literatura (Letizia y
Lillo, 2019; Óskarsdóttir y Bravo, 2021). La sincronización no aparece en esos
trabajos.}
\end{preguntas}

\diapo{4}{Seis tareas, un solo modelo}{1:15}{5:30}{Cada pieza se valida antes de
agregar la siguiente.}
\begin{decir}
El proyecto sigue la propuesta del semillero. Las tareas uno a cinco construyen el
modelo: primero la red y el déficit, luego la redistribución de liquidez, luego la
sincronización, luego los impagos en cascada y luego la predicción con aprendizaje
automático. La tarea seis, en rojo, aplica todo a la cartera de un banco.

La regla de construcción es estricta: cada tarea agrega un término a la ecuación, y si
ese término se apaga, el modelo vuelve exactamente al de la tarea anterior. Así cada
pieza se puede probar por separado, y si algo falla sabemos exactamente en qué paso
ocurrió. No presentaré las seis tareas con el mismo detalle: mostraré lo esencial de las
cinco primeras y me concentraré en la sexta, que responde la pregunta.
\transicion{Veamos la ecuación completa.}
\end{decir}
\begin{terminos}
\begin{description}
\item[Aprendizaje automático (ML).] Modelos que aprenden patrones a partir de datos para
hacer predicciones.
\item[Cascada.] Impagos en cadena: el impago de uno lleva a otros a no poder pagar.
\end{description}
\end{terminos}

% ===================================================================
\section{El modelo}

\diapo{5}{La ecuación del déficit}{1:45}{7:15}{Cinco efectos en una ecuación.}
\begin{decir}
Esta ecuación dice cómo cambia en el tiempo el déficit de cada entidad. La letra $x_i$
es el déficit de la entidad $i$: cuánta liquidez le falta. La derivada $dx_i/dt$ es qué
tan rápido sube o baja.

Cada término, de izquierda a derecha, es un efecto.

En verde, menos alfa por $x_i$: la corrección. Cada entidad reduce su déficit por sí
misma, por ejemplo con ahorro. Alfa es qué tan rápido lo hace.

En rojo, beta por la suma de $W_{ij}$ por $x_j$: el contagio. $W_{ij}$ es cuánto le paga
$j$ a $i$. Si quienes le pagan a $i$ tienen déficit, le pagan menos y su déficit sube.
Beta es la intensidad del contagio.

En azul, la redistribución: la liquidez se traslada entre entidades sin crearse ni
destruirse. $D_0$ es su velocidad y $L$ es una matriz llamada Laplaciano.

En morado, el desfase: gamma por $k_i$ por uno menos $r_i$. $k_i$ es el ingreso total de
$i$ y $r_i$ dice si ese ingreso llega a tiempo. Lo explico en la siguiente diapositiva.

Y al final, eta: los choques externos.

Cuando el déficit supera el colchón de la entidad, la entidad deja de pagar. Eso es un
impago.
\transicion{Falta definir $r_i$, que viene del ritmo de pagos.}
\end{decir}
\begin{terminos}
\begin{description}
\item[$x_i$.] Déficit de la entidad $i$. Tareas 1–5: en pagos medios por ciclo. Tarea 6:
fracción del ingreso mensual.
\item[$t$.] Tiempo: ciclos de pago (un día hábil) en las tareas 1–5; meses en la tarea 6.
\item[$\alpha$.] Tasa de corrección [1/tiempo]. $1/\alpha$ es el tiempo típico en que se
corrige un déficit.
\item[$\beta$.] Intensidad del contagio de déficit de pagador a receptor.
\item[$W_{ij}$.] Pago medio de $j$ hacia $i$ en cada ciclo. Las filas son ingresos; las
columnas, pagos.
\item[$\sum_jW_{ij}x_j$.] Déficit de todos los pagadores de $i$, pesado por cuánto le
paga cada uno.
\item[$D_0$, $L$.] Velocidad de redistribución y Laplaciano de flujo, $L=D_\text{out}-W$,
con $D_\text{out}$ la matriz diagonal de pagos totales. Sus columnas suman cero, así que
conserva el déficit total.
\item[$\gamma$.] Cuánto déficit genera el desfase [1/tiempo].
\item[$k_i$.] Ingreso de $i$ por la red, $\sum_jW_{ij}$.
\item[$r_i$.] Coherencia local, entre $-1$ y $1$ (diapositiva 6).
\item[$\eta_i(t)$.] Choques externos: retiros de liquidez o ruido aleatorio.
\item[$x_i^{(c)}$.] Colchón o umbral de impago de la entidad $i$.
\end{description}
\end{terminos}
\begin{preguntas}
\pr{¿Por qué lineal?}{Es lo más simple que conserva el contagio y permite resultados
exactos, como el umbral $R=1$. La no linealidad entra por el umbral y por el ritmo. Es
una limitación declarada.}
\end{preguntas}

\diapo{6}{La ecuación del ritmo de pagos}{1:45}{9:00}{El término $\gamma k_i(1-r_i)$ une
el ritmo con la liquidez.}
\begin{decir}
Cada entidad paga en ciclos, como un reloj que da vueltas. Theta sub $i$ es la posición
de esa manecilla: en qué punto de su ciclo de pagos está.

La ecuación de la izquierda es el modelo de Kuramoto, el modelo clásico de
sincronización en física. Omega sub $i$ es el ritmo propio de la entidad, el que tendría
sola. El segundo término la acelera si quienes le pagan van adelante y la frena si van
atrás. $K$ es qué tan fuerte se ajusta. Es lo mismo que pasa con un público que aplaude
y termina aplaudiendo al unísono.

La fórmula de la derecha es la coherencia, $r_i$. Promedia el coseno de la diferencia de
fase con cada pagador, pesado por lo que paga. Si el ingreso llega a tiempo, $r_i$ vale
uno. Si llega en el peor momento, vale menos uno. En el dibujo, tres pagadores están en
fase y uno desfasado.

Ahora el término morado de la ecuación anterior tiene sentido: uno menos $r_i$ es cero
si todo llega a tiempo, y es grande si llega desfasado. Ese término convierte el desfase
en déficit.
\transicion{Antes de la aplicación a crédito, tres resultados del modelo.}
\end{decir}
\begin{terminos}
\begin{description}
\item[$\theta_i$.] Fase: posición en el ciclo de pagos, en radianes, de $0$ a $2\pi$.
\item[$\omega_i$.] Frecuencia natural: ritmo propio de la entidad [rad/tiempo]. Se sortea
de una distribución de Lorentz con ancho $\Delta$.
\item[$K$.] Fuerza del acoplamiento de fase.
\item[$\langle k\rangle$.] Ingreso medio de la red. Normalizar por él, y no por el número
de entidades, evita que el acople desaparezca en redes grandes y dispersas (Rodrigues et
al., 2016).
\item[$\sin(\theta_j-\theta_i)$.] Positivo si el pagador va adelante: acelera a $i$.
\item[$r_i$.] Coherencia local $\frac1{k_i}\sum_jW_{ij}\cos(\theta_j-\theta_i)$.
\item[Modelo de Kuramoto.] Modelo de osciladores acoplados propuesto por Yoshiki
Kuramoto en 1975.
\end{description}
\end{terminos}

% ===================================================================
\section{Lo que muestra el modelo}

\diapo{7}{La red amortigua los choques solo si $R<1$}{1:15}{10:15}{Con $R=0.9$ el déficit
se apaga; con $R=1.1$ crece.}
\begin{decir}
El primer resultado responde cuándo la red amortigua un choque y cuándo lo amplifica.

La respuesta es un solo número, $R$: beta por rho de $W$ dividido por alfa. Arriba va el
contagio; abajo, la corrección. Rho de $W$ es el radio espectral de la red, su valor
propio más grande, que mide qué tan concentrados están los pagos.

Funciona como el número reproductivo de una epidemia. Las gráficas muestran el déficit
de cada entidad después de un choque. El eje horizontal es el tiempo y el vertical el
déficit, en escala logarítmica. Cada línea de color es una entidad y la línea negra es
el total. Con $R$ igual a 0.9, a
la izquierda, el total baja. Con $R$ igual a 1.1, a la derecha, crece sin límite. Esto
lo comprobamos contra la solución exacta, con error menor que una millonésima.
\transicion{El segundo resultado es sobre el ritmo.}
\end{decir}
\begin{terminos}
\begin{description}
\item[$\rho(W)$ (radio espectral).] Mayor valor propio de $W$ en módulo. Por el teorema de
Perron–Frobenius, en una matriz de pagos positivos es real y positivo.
\item[Valor propio.] Número $\lambda$ tal que $Wv=\lambda v$ para algún vector $v$.
\item[Número reproductivo ($R_0$).] En epidemiología, contagios que produce en promedio
un infectado; la epidemia crece si es mayor que 1.
\item[Solución exacta.] Sin umbrales, el sistema es lineal y se resuelve como
$x(t)=e^{Mt}x_0$.
\end{description}
\end{terminos}

\diapo{8}{Los ciclos de pago se sincronizan como predice la teoría}{1:00}{11:15}{El
código reproduce un resultado exacto de la física.}
\begin{decir}
El segundo resultado valida la parte del ritmo. La gráfica muestra la sincronía
global, $r$, en el eje vertical, contra el acoplamiento, $K$, en el horizontal. Cada
color es un tipo de red. Cuando $r$ vale
cero, cada entidad paga a su ritmo; cuando vale uno, todas pagan a la vez. Al aumentar
$K$, la red se sincroniza.

La línea punteada es la teoría para una red donde todos están conectados con todos. Para
$K$ igual a 2, la teoría da 0.7071 y la simulación 0.7076. La diferencia está en la
cuarta cifra decimal. Las redes con menos conexiones se sincronizan de forma más gradual,
pero siguen la misma tendencia.
\transicion{El tercer resultado es sobre las cascadas.}
\end{decir}
\begin{terminos}
\begin{description}
\item[$r$ global.] $r=\bigl|\frac1N\sum_je^{i\theta_j}\bigr|$, entre 0 y 1.
\item[$\Delta$.] Dispersión de los ritmos propios. Teoría: $r=\sqrt{1-2\Delta/K}$ para
$K>2\Delta$.
\item[Topologías.] Erdős–Rényi (conexiones al azar), Barabási–Albert (pocos nodos muy
conectados) y núcleo-periferia (núcleo denso, como en los sistemas de pagos reales).
\end{description}
\end{terminos}

\diapo{9}{Las cascadas son máximas con conectividad intermedia}{1:00}{12:15}{Más
conexiones no siempre significa más riesgo.}
\begin{decir}
Aquí ya hay impagos. Chocamos una entidad y medimos qué fracción de la red termina en
impago, según cuántas conexiones tiene cada entidad en promedio.

Con pocas conexiones, el choque no tiene por dónde propagarse. Con muchas, cada entidad
recibe de muchos pagadores y la pérdida de uno se diluye. En el medio, las cascadas son
máximas: en la gráfica, la fracción en impago sube rápido, llega a su máximo con un grado
medio cercano a cinco y medio, y luego cae. Eso se conoce como ventana de contagio. Nuestro modelo, en azul, la reproduce,
y se diferencia de dos modelos clásicos, DebtRank y Watts, que no tienen tiempo ni
recuperación.
\transicion{Simular todo esto es lento; por eso entrenamos modelos predictivos.}
\end{decir}
\begin{terminos}
\begin{description}
\item[Grado medio.] Número medio de contrapartes por entidad.
\item[Colchón.] $x_i^{(c)}=c\,\mathrm{out}_i\,\varepsilon_i$: múltiplo $c$ de lo que la
entidad paga por ciclo, con variación aleatoria $\varepsilon_i$.
\item[DebtRank.] Medida de propagación de estrés financiero (Battiston et al., 2012).
\item[Modelo de Watts.] Un nodo cae si la fracción de vecinos caídos supera un umbral
(Watts, 2002).
\end{description}
\end{terminos}

\diapo{10}{El aprendizaje automático predice el colapso en milisegundos}{1:15}{13:30}{El
modelo físico genera los datos; el aprendizaje automático los resume.}
\begin{decir}
Generamos 20\,000 simulaciones sobre 300 redes distintas y entrenamos modelos para
predecir si la red colapsa, es decir, si cae el 30~\% o más.

El mejor, XGBoost, alcanza un AUC de 0.98. Más adelante explico el AUC; por ahora, 1 es
perfecto. Una simulación tarda 0.08 segundos y una predicción 0.2 milisegundos.

Esa diferencia de velocidad
es lo que permite que la aplicación muestre el riesgo de todas las entidades al instante.

Para evaluar usamos siempre redes que el modelo no vio durante el entrenamiento. La
gráfica muestra el AUC por grupos: a la izquierda por rangos de $R$ y a la derecha por
tipo de red. XGBoost, en naranja, se mantiene alto en todos los grupos.
\transicion{Con el modelo construido y validado, pasamos a la aplicación: el crédito.}
\end{decir}
\begin{terminos}
\begin{description}
\item[XGBoost.] Conjunto de árboles de decisión entrenados uno tras otro, cada uno
corrigiendo los errores del anterior.
\item[Regresión logística.] Modelo lineal para probabilidades; es la referencia.
\item[AUC.] Área bajo la curva ROC (diapositiva 15).
\item[Validación por grupos (GroupKFold).] Ninguna red aparece a la vez en entrenamiento
y en evaluación.
\end{description}
\end{terminos}
\begin{preguntas}
\pr{¿También usaron redes neuronales?}{Sí: una red neuronal sobre grafos (GCN) para
predecir el tamaño de la cascada de cada entidad, y una red informada por la física
(PINN) para estimar parámetros a partir de trayectorias. Están en la diapositiva de
respaldo 5.}
\end{preguntas}

% ===================================================================
\section{Aplicación al riesgo de crédito}

\diapo{11}{Una cartera bancaria en tres capas}{1:15}{14:45}{Se simula una cartera para
comprobar que el método encuentra lo que sabemos que está.}
\begin{decir}
Ahora los nodos son clientes de un banco: personas y empresas. Se relacionan de tres
formas. La nómina: la empresa paga a su empleado cada quince o treinta días. Los
proveedores: una empresa paga a otra sus facturas. Y las transferencias entre personas,
que son irregulares.

No existen datos públicos de la red de pagos de los clientes de un banco. Por eso
generamos una cartera: 2000 personas y 100 empresas, con 55~\% de informalidad, una cifra
del orden de la que reporta el DANE. Le ponemos un mecanismo conocido y verificamos que
el método lo encuentra. Es como calibrar un instrumento: primero se mide algo cuyo valor
se conoce.

La figura es una versión reducida para ver las capas: los cuadrados son empresas, los
puntos son personas, y el color de cada línea indica la capa.
\transicion{El banco, sin embargo, no ve el modelo; solo ve los pagos.}
\end{decir}
\begin{terminos}
\begin{description}
\item[Cartera.] Conjunto de clientes con crédito en un banco.
\item[Red multicapa.] Red con varios tipos de relación entre los mismos nodos; cada tipo
es una capa.
\item[Informalidad.] Personas sin empleador dentro del banco. Sus ingresos llegan con
ritmos irregulares.
\item[Datos sintéticos.] Generados por simulación con reglas conocidas.
\item[PD verdadera.] Cada cliente tiene una probabilidad de impago oculta, calculada con
sus variables clásicas y un componente que el banco no ve. Los modelos tienen prohibido
usarla: el código se detiene si lo intentan.
\end{description}
\end{terminos}
\begin{preguntas}
\pr{¿Qué vale un resultado con datos sintéticos?}{Demuestra que, si el mecanismo
existe, el método lo detecta con la información que tiene un banco, y que no detecta
nada si no existe. No demuestra que exista en un banco real.}
\end{preguntas}

\diapo{12}{El banco solo ve transacciones; el ritmo se estima desde ellas}{1:15}{16:00}{La
coherencia estimada coincide con la real.}
\begin{decir}
Este es el paso clave para que el método sea aplicable.

Primero, el modelo genera los pagos tal como los registraría un banco: fecha, quién paga,
quién recibe y cuánto. Esa tabla es lo único que usamos.

Segundo, con esos pagos estimamos la fase de cada cliente. La idea es sencilla: entre un
pago y el siguiente, el cliente completó una vuelta de su ciclo. Entonces, a mitad de
camino entre dos pagos, está a media vuelta.

Tercero, con las fases estimadas calculamos la coherencia de cada cliente con sus
pagadores.

La figura compara lo real con lo estimado para un empleado formal y para un informal. La
correlación entre la coherencia real y la estimada es mayor que 0.8.

Usamos seis meses de transacciones para observar y predecimos si el cliente cae en los
doce meses siguientes.
\transicion{Con eso construimos las variables.}
\end{decir}
\begin{terminos}
\begin{description}
\item[Registro de transacciones.] Tabla con fecha, pagador, receptor, tipo de pago y
monto.
\item[Fase por eventos.] Pikovsky et al.\ (2001):
$\hat\theta(t)=2\pi\bigl(k+\frac{t-t_k}{t_{k+1}-t_k}\bigr)$, donde $t_k$ es la fecha del
$k$-ésimo pago.
\item[$\hat r_i$.] Coherencia estimada con las fases $\hat\theta$ y la red estimada
$\hat W$. El acento circunflejo indica ``estimado''.
\item[Ventanas.] Observación: meses 0 a 6. Predicción: meses 6 a 18.
\item[Fuga de información.] Usar al entrenar algo que no se tendría al predecir. El código
rechaza pagos posteriores al mes 6.
\end{description}
\end{terminos}

\diapo{13}{Tres niveles de información}{1:00}{17:00}{El nivel 3 es el aporte del
proyecto.}
\begin{decir}
Comparamos tres conjuntos de variables, y cada uno contiene al anterior.

El nivel uno es el score clásico: ingreso, deuda, edad, antigüedad, moras y si el empleo
es formal.

El nivel dos agrega la red estática: de quién recibe el cliente, cuánto, qué tan
concentrado está su ingreso en pocos pagadores y qué tan sano está su empleador.

El nivel tres agrega la dinámica: la coherencia, el desfase con sus pagadores, la
diferencia de ritmo, qué tan regulares son sus ciclos y si algún pagador dejó de pagar.

Los tres se evalúan con el mismo modelo, los mismos clientes y la misma validación.
Así, cualquier diferencia de desempeño se debe solo a la información agregada.
\transicion{Y los evaluamos en dos mundos.}
\end{decir}
\begin{terminos}
\begin{description}
\item[Concentración (HHI).] Índice de Herfindahl–Hirschman: suma de las participaciones
de cada pagador al cuadrado. Vale 1 si todo viene de uno solo.
\item[Salud del empleador.] Tamaño, deuda y moras de la empresa que paga la nómina.
\item[Desfase.] Diferencia media de fase con los pagadores.
\item[Regularidad.] Variación de los intervalos entre pagos (coeficiente de variación).
\end{description}
\end{terminos}

\diapo{14}{El experimento compara dos mundos idénticos}{1:00}{18:00}{Si la dinámica
mejorara en el mundo nulo, el método estaría engañando.}
\begin{decir}
Este es el diseño que hace creíble el resultado.

Simulamos dos mundos con los mismos clientes, la misma red y los mismos números
aleatorios. En el mundo con red, el desfase y el contagio generan déficit. En el mundo
nulo, los pagos ocurren igual, pero no generan déficit.

En el mundo con red, la dinámica debería ayudar a predecir. En el nulo no debería. Si
ayudara en el nulo, significaría que el método encuentra patrones que no existen, por
ejemplo por una fuga de información.

Es un control negativo, como el placebo en un ensayo clínico.

Usamos 20 carteras por mundo, unos 40\,000 clientes, y validación cruzada: el modelo
siempre se evalúa en carteras que no vio al entrenar.
\transicion{Veamos qué pasó.}
\end{decir}
\begin{terminos}
\begin{description}
\item[Mundo nulo.] $\beta=\gamma=\lambda=0$: sin contagio, sin desfase y sin pérdida por
pagadores caídos.
\item[Mismos números aleatorios.] La semilla de cada cartera genera los mismos sorteos en
ambos mundos.
\item[Validación cruzada.] Se divide en 5 partes; se entrena con 4 y se evalúa en la
quinta, rotando.
\end{description}
\end{terminos}

% ===================================================================
\section{Resultados}

\diapo{15}{Cada nivel de información mejora la predicción}{1:15}{19:15}{La dinámica
agrega información que la red estática no tiene.}
\begin{decir}
Primero, qué es el AUC. Tomemos al azar un cliente que cayó y uno que no cayó. El AUC es
la probabilidad de que el modelo le asigne más riesgo al que cayó. Si vale 0.5, el
modelo adivina al azar; si vale 1, acierta siempre.

En el mundo con red, el score clásico logra 0.769. Al agregar la red estática sube a
0.784. Al agregar la dinámica sube a 0.795. Con XGBoost el resultado final es el mismo.
Todas estas cifras se miden en carteras que el modelo no vio.

Puede parecer un cambio pequeño, pero hay que leerlo en dos pasos. El paso de la red
estática ya estaba documentado en la literatura. El paso de la dinámica es el nuevo, y se
suma encima del anterior: la dinámica no repite lo que ya decía la red estática.
\transicion{¿Es real esa mejora, o es ruido?}
\end{decir}
\begin{terminos}
\begin{description}
\item[AUC.] Área bajo la curva ROC. La curva ROC grafica aciertos contra falsas alarmas
para todos los umbrales de decisión.
\item[Fuera de muestra.] Medido en datos que el modelo no usó para aprender.
\item[Gini.] Métrica usada por los bancos: $2\,\mathrm{AUC}-1$.
\end{description}
\end{terminos}
\begin{preguntas}
\pr{¿Por qué el AUC del mundo nulo es mayor (0.81)?}{En el nulo el impago depende casi
solo de la probabilidad propia de cada cliente, que es lo que mide el score. Con red,
parte del impago viene de la red y el score no lo ve.}
\end{preguntas}

\diapo{16}{La mejora es real: aparece con red y desaparece en el control}{1:15}{20:30}{Con
red, positiva en 5 de 5 particiones; en el nulo, cero.}
\begin{decir}
Cómo leer la gráfica: cada fila es una comparación. El punto es la mejora promedio de la
dinámica sobre la red estática y la barra es su intervalo de confianza del 95~\%. La
línea punteada vertical es el cero: si una barra la cruza, no podemos afirmar que haya
mejora.

En el mundo con red, en rojo, la mejora es de 0.011 con regresión logística y 0.014 con
XGBoost. Los intervalos no tocan el cero y la mejora es positiva en las cinco
particiones.

En el mundo nulo, en gris, la mejora es cero. Esto es lo más importante de la
presentación: con las mismas variables y el mismo método, cuando el mecanismo no existe,
no aparece ninguna mejora. Eso descarta que el resultado sea un artefacto.
\transicion{Y podemos ir más allá: medir cómo cambia la mejora con el mecanismo.}
\end{decir}
\begin{terminos}
\begin{description}
\item[$\Delta$AUC.] AUC del nivel 3 menos AUC del nivel 2, en la misma partición.
\item[Intervalo de confianza 95~\%.] Rango que contiene la mejora verdadera con 95~\% de
confianza. Se calcula con la distribución $t$ de Student sobre las 5 particiones.
\item[Partición.] Cada una de las 5 divisiones de la validación cruzada.
\end{description}
\end{terminos}
\begin{preguntas}
\pr{¿Una mejora de 0.01 importa?}{En crédito sí: son 0.022 puntos de Gini sobre un
modelo que ya usa la red. Además, el valor depende de parámetros ilustrativos; lo que se
establece es que la mejora existe, es estable y desaparece sin el mecanismo.}
\pr{¿Más variables no mejoran siempre?}{No fuera de muestra. El mundo nulo lo prueba: las
mismas variables, mejora cero.}
\end{preguntas}

\diapo{17}{La mejora crece con la intensidad del desfase}{1:00}{21:30}{La señal viene de
la desincronización, no del contagio.}
\begin{decir}
Este experimento responde una pregunta que surge naturalmente: ¿la mejora viene del
desfase o del contagio? Para saberlo variamos solo gamma, el parámetro que convierte el
desfase en déficit. Todo lo demás queda igual, incluido el contagio.

En la gráfica, el eje horizontal es gamma y el vertical la mejora de la dinámica sobre la
red estática.

Con gamma igual a cero, la dinámica no aporta nada, aunque el contagio siga activo. Al
subir gamma, la mejora crece: 0.001, 0.008 y 0.029.

Esto nos dice de dónde viene la señal: de la desincronización, no del contagio. Y es una
relación de dosis y respuesta: más mecanismo, más señal.
\transicion{¿Significa esto que la red reemplaza al score? No.}
\end{decir}
\begin{terminos}
\begin{description}
\item[$\gamma$.] Déficit generado por el desfase [1/mes].
\item[Dosis-respuesta.] El efecto crece con la intensidad de la causa; es un criterio
clásico para atribuir un efecto a una causa.
\end{description}
\end{terminos}

\diapo{18}{El score clásico sigue siendo la base}{0:45}{22:15}{El modelo híbrido
complementa el score; no lo reemplaza.}
\begin{decir}
La gráfica muestra cuánto pesa cada variable en la predicción, con un método llamado
SHAP.

Cada barra es una variable; mientras más larga, más influye en la predicción. Las
variables clásicas, en gris, son las más importantes: la deuda, el ingreso y las
moras. Las de red estática, en azul, y las dinámicas, en rojo, agregan información. Entre
las dinámicas pesan más la variabilidad de la coherencia, su nivel y el desfase.

La conclusión práctica es que el modelo híbrido complementa al score, no lo reemplaza.
\transicion{Un ejemplo de uso.}
\end{decir}
\begin{terminos}
\begin{description}
\item[SHAP.] \emph{SHapley Additive exPlanations} (Lundberg y Lee, 2017). Reparte cada
predicción entre las variables usando los valores de Shapley de la teoría de juegos.
\item[Importancia media $|\mathrm{SHAP}|$.] Promedio del tamaño de la contribución de cada
variable.
\item[Modelo híbrido.] Modelo con los tres niveles de variables.
\end{description}
\end{terminos}

\diapo{19}{Uso práctico: la quiebra de un empleador}{1:15}{23:30}{Sirve para pruebas de
estrés por concentración en empleadores.}
\begin{decir}
Además de predecir, el modelo permite simular escenarios. Simulamos que quiebra la
empresa con más empleados en el banco, 74 personas, en el primer mes. La gráfica muestra
el porcentaje de esos empleados en impago a lo largo de siete meses: en rojo con
quiebra, en línea punteada sin ella, y en gris el mundo nulo.

En el mundo con red, el 12~\% de sus empleados cae de inmediato, porque pierde su salario.
A los seis meses, el 13~\% está en impago, frente al 7~\% si la empresa no hubiera
quebrado. En el mundo nulo no hay efecto, como corresponde.

Para un banco, esto es una prueba de estrés: cuánto riesgo concentra en los empleados de
una misma empresa.
\transicion{Para cerrar los resultados, qué podemos afirmar y qué no.}
\end{decir}
\begin{terminos}
\begin{description}
\item[Prueba de estrés.] Simulación de un escenario adverso para medir su impacto.
\end{description}
\end{terminos}

\diapo{20}{Qué demuestra y qué no}{1:00}{24:30}{El siguiente paso es una cartera real
anonimizada.}
\begin{decir}
El trabajo demuestra tres cosas. Si el mecanismo existe, se detecta con las transacciones
que ya tiene un banco. Aporta información que la red estática no tiene. Y desaparece
cuando el mecanismo no está, y crece cuando es más intenso.

No demuestra que el mecanismo exista en un banco real, porque los datos son simulados.
No da el tamaño real del efecto, porque los parámetros son ilustrativos. Y supone que el
banco ve toda la red de sus clientes. El paso siguiente es probarlo con una cartera real
anonimizada.
\transicion{Termino con el software y las conclusiones.}
\end{decir}
\begin{preguntas}
\pr{¿Qué haría falta para usarlo en un banco?}{Una cartera real anonimizada, bajo la Ley
1581 de 2012 de protección de datos. El código ya recibe transacciones con esa
estructura.}
\end{preguntas}

% ===================================================================
\section{Cierre}

\diapo{21}{Software reproducible y aplicación web}{0:45}{25:15}{Todo resultado se puede
regenerar y explorar.}
\begin{decir}
Todo está implementado en Python, con una aplicación web de nueve páginas, como la
que se ve a la derecha, donde se puede construir una red, aplicar choques y ver la
predicción de cada cliente. Los
parámetros están en archivos de configuración, las semillas aleatorias son fijas y hay
127 pruebas automáticas. Con Docker, la aplicación funciona igual en cualquier
computador, y el experimento completo de crédito se repite con un solo comando en unos
20 segundos.
\end{decir}
\begin{terminos}
\begin{description}
\item[Streamlit.] Biblioteca de Python para aplicaciones web interactivas.
\item[YAML.] Formato de texto para archivos de configuración.
\item[Docker.] Empaqueta el programa con todo lo que necesita para que corra igual en
cualquier equipo.
\item[Prueba automática.] Código que verifica un resultado; se ejecuta con
\texttt{pytest}.
\end{description}
\end{terminos}

\diapo{22}{Conclusiones}{1:15}{26:30}{No solo importa cuánto recibe un cliente, también
cuándo.}
\begin{decir}
Cuatro conclusiones.

Primera: un solo modelo une contagio, redistribución, sincronización y cascadas, y cada
parte se comprobó contra resultados exactos.

Segunda: el número $R$ dice si la red amortigua o amplifica un choque.

Tercera: el aprendizaje automático predice el colapso con AUC de 0.98 en milisegundos.

Cuarta, y es la principal: la sincronización de los pagos, medida solo con
transacciones, mejora la predicción del impago más allá del score y de la red estática,
y no lo hace cuando el mecanismo no existe.

Para la física, el proyecto muestra que las herramientas de sincronización y redes
complejas tienen una aplicación directa en finanzas. Para un banco, sugiere una fuente
de información que ya tiene en sus registros y que hoy no usa.

En resumen: no solo importa cuánto recibe un cliente, también cuándo.
\end{decir}

\diapo{23}{Trabajo futuro}{0:30}{27:00}{Validación con datos reales.}
\begin{decir}
Como trabajo futuro: probar el método con datos reales, incluir las fechas de cuota como
un ritmo impuesto, permitir que el déficit cambie el ritmo de pago, y estudiar qué pasa
cuando un banco solo ve a sus clientes, frente al operador de Bre-B, que ve todo.
\end{decir}

\diapo{Cierre}{Gracias}{0:15}{27:15}{Abrir preguntas.}
\begin{decir}
Muchas gracias. Quedo atento a sus preguntas.
\end{decir}
